\documentclass[11pt]{article}
\usepackage[T1]{fontenc}
\usepackage{textcomp}
\usepackage[margin=0.75in]{geometry}
\usepackage{amsmath,amssymb,amsthm}
\usepackage{array,booktabs}
\usepackage{hyperref}
\setlength{\parindent}{0pt}
\newtheorem{theorem}{Theorem}
\theoremstyle{definition}
\newtheorem{definition}{Definition}
\title{Flow Proof Artifact: Comb-choose-two}
\author{Generated by \texttt{flow doc proof}}
\date{Flow Proof Book}
\begin{document}
\maketitle
Choosing two items relates to choosing one by Pascal.\\[0.5em]
\textbf{Source.} graham-knuth-patashnik — *Concrete Mathematics*\\[0.5em]
\bigskip
\noindent{\large \textbf{Derived fact 1.} \textit{choose(n, 2) = choose(n, 1) + choose(n - 1, 1) for n > 1} \hfill Comb \cdot choose \cdot choosing two via Pascal}\par
\smallskip\noindent\textit{Coordinate: Comb \cdot choose \cdot choosing two via Pascal}\par
\smallskip\noindent\textit{Source: graham-knuth-patashnik}\par
\smallskip\noindent\textit{Built on: choosing one gives the count, for choose on Comb, pascal recurrence holds, for choose on Comb}\par
\medskip
\noindent\textbf{Goal.} choose(n, 2) = choose(n, 1) + choose(n - 1, 1) for n > 1\par
\noindent$\displaystyle \forall n \in \mathbb{N}\quad choose(n, 2) = choose(n, 1) + choose(pred(n), 1)$\par
\medskip
\noindent
\begin{tabular}{@{} >{\raggedright\arraybackslash}p{0.06\textwidth} >{\raggedright\arraybackslash}p{0.47\textwidth} >{\raggedleft\arraybackslash}p{0.41\textwidth} @{}}
\textbf{\#} & \textbf{Proof} & \textbf{Mathematics} \\
\midrule
\textbf{1.} & We prove that choosing two via Pascal for choose on Comb. & \\[0.45em]
\textbf{2.} & We split into exhaustive cases — the claim must hold in each one. & \\[0.45em]
\textbf{3.} & Case 1 (see step 2): suppose n is zero. & \\[0.45em]
\textbf{4.} & We invoke the derived fact governing choose on Comb: choosing one gives the count, for choose on Comb (instantiated for 0). & \\[0.45em]
\textbf{5.} & From step 3 and step 4, this implies choose(n, 2) equals choose(n, 1) plus choose(pred(n), 1) in this case. & $\displaystyle choose(n, 2) = choose(n, 1) + choose(pred(n), 1)$ \\[0.45em]
\textbf{6.} & Case 2 (see step 2): neither disjunct holds. & \\[0.45em]
\textbf{7.} & Let k = 2. & \\[0.45em]
\textbf{8.} & We invoke the definitional clause governing choose on Comb: pascal recurrence holds, for choose on Comb (instantiated for n, k). & \\[0.45em]
\textbf{9.} & We invoke the derived fact governing choose on Comb: choosing one gives the count, for choose on Comb (instantiated for n). & \\[0.45em]
\textbf{10.} & From step 6, step 7, step 8, and step 9, this implies choose(n, 2) equals choose(n, 1) plus choose(pred(n), 1). Together with the other cases (step 3 and step 6), the goal is discharged. Hence proven. & $\displaystyle choose(n, 2) = choose(n, 1) + choose(pred(n), 1)$ \\[0.45em]
\bottomrule
\end{tabular}
\label{thm:Comb-choose-choosing_two_via_pascal}
\par\medskip\hrule
\end{document}
